\chapter{William Thurston}

\addcontentsline{toc}{chapter}{Problem History}
\section*{Problem History}

\begin{prereqs}
\noindent \textbf{Background}
\begin{itemize}
\item Topology of 3-manifolds
\item Hyperbolic geometry and geometric structures
\item Group theory and geometrization
\end{itemize}
\end{prereqs}

\begin{objectives}
\noindent \textbf{Learning Objectives}
\begin{itemize}
\item Understand Thurston's geometrization conjecture and its eight model geometries
\item Learn how hyperbolic geometry characterizes the complexity of 3-manifolds
\item Appreciate the unifying vision of geometrization and its implications for topology
\end{itemize}
\end{objectives}

Thurston \cite{thurston1997} transformed 3-manifold study from topological to geometric. His
geometrization conjecture proposes that every closed 3-manifold decomposes
along spheres and tori into pieces each carrying one of eight model
geometries: $\mathbb{E}^3$, $S^3$, $\mathbb{H}^3$, $S^2 \times \mathbb{R}$,
$\mathbb{H}^2 \times \mathbb{R}$, $\widetilde{SL(2,\mathbb{R})}$, Nil,
and Sol.

The prime decomposition theorem (Kneser--Milnor) splits along spheres
into primes. Thurston's conjecture further splits along essential tori
via the JSJ decomposition into geometric pieces.

A geometry on a 3-manifold $M$ is a Riemannian metric whose isometry
group acts transitively and whose universal cover is one of the eight
model spaces. A $(G, X)$-structure---an atlas of charts into $X$ with
transition maps in $G$---formalizes this. The developing map sends the
universal cover to $X$, with holonomy representation
$\rho: \pi_1(M) \to G$ determining the structure.

Hyperbolic 3-manifolds are ``generic'': most 3-manifolds admit a
hyperbolic structure. Mostow rigidity ensures uniqueness for
finite-volume hyperbolic structures, making volume a topological
invariant.

\addcontentsline{toc}{chapter}{How Was It Solved}
\section*{How Was It Solved}

Thurston proved geometrization for Haken manifolds (containing essential
surfaces). His hyperbolic Dehn surgery theorem---that filling cusps of a
hyperbolic 3-manifold yields hyperbolic structures for all but finitely
many coefficients---connects cusped to closed manifolds.

The ending lamination conjecture (proved by Brock--Canary--Minsky, 2012)
provides a complete description of Kleinian group deformation spaces: a
hyperbolic 3-manifold is determined by its end invariants (conformal
structure at infinity or ending laminations).

For Haken manifolds, the proof proceeds by induction on complexity. Normal
surface theory provides a combinatorial framework for analyzing surface
intersections. Essential surfaces allow decomposition into simpler pieces
or identification of one of the eight geometries.

For hyperbolic pieces, quasi-Fuchsian group theory and Bers simultaneous
uniformization apply: the quasi-Fuchsian deformation space of a surface
group is parameterized by conformal structure pairs, with ending
laminations providing complete descriptions.

\addcontentsline{toc}{chapter}{Impact}
\section*{Impact}

Geometrization provided a unifying framework for 3-manifold topology,
encompassing the Poincar\'{e} conjecture, prime decomposition, and JSJ
decomposition. Hyperbolic volume---a topological invariant by Mostow
rigidity---connects hyperbolic geometry with combinatorial topology.

The theory of $(G, X)$-structures has been applied to higher dimensions.
The ending lamination theorem deepened the topology-holonomy
correspondence. In physics, hyperbolic volume appears in Chern-Simons
invariant perturbative expansions; the Volume Conjecture links the
colored Jones polynomial to hyperbolic volume.

\addcontentsline{toc}{chapter}{Modern Views}
\section*{Modern Views}

Perelman's proof of geometrization (2003) using Ricci flow with surgery
confirms Thurston's vision and resolves the Poincar\'{e} conjecture. The
virtual Haken and virtual fibering conjectures follow (Agol, 2012).

Active research continues on exceptional Dehn surgeries, small-volume
classification, and volume-invariant relationships. Thurston's geometric
perspective influenced geometric group theory (asymptotic cones,
hyperbolic and CAT(0) groups), with the interaction between 3-manifold
topology and cubulated groups being especially productive.

Thurston received the Fields Medal in 1982 for 3-manifold topology,
geometrization, and hyperbolic structures.

\section*{Summary}
Thurston's geometrization conjecture proposes that every closed 3-manifold
decomposes into pieces admitting one of eight model geometries. His work
on hyperbolic manifolds, Kleinian groups, and deformation theory provided
the tools. Perelman's Ricci flow proof confirms the vision and has
transformed 3-manifold topology.

\addcontentsline{toc}{chapter}{Exercises}
